% STATUS: PHASE 1B ASSESSMENT; PROJECT THEOREMS REMAIN UNPROVED
\section{Revised novelty matrix: eight adversarial decisions}
\label{sec:revised-matrix}
This matrix replaces the original matrix only for current planning; the
original file remains preserved. Exact source assumptions, locations and
access limitations are recorded in Section~\ref{sec:adversarial-audit}.
RED means known, routine or unsafe as novelty; ORANGE means substantial
overlap requiring strengthening; YELLOW means a potentially new but modest
specialization; GREEN would mean an apparently substantive new result;
GRAY means undetermined. No GREEN finding is warranted. ``Not located''
is never a positive novelty argument.

\begin{longtable}{@{}>{\raggedright\arraybackslash}p{0.07\textwidth}>{\raggedright\arraybackslash}p{0.28\textwidth}>{\raggedright\arraybackslash}p{0.13\textwidth}>{\raggedright\arraybackslash}p{0.44\textwidth}@{}}
\toprule
ID & Target and owners & Verdict & Strongest adverse evidence / present use\\
\midrule\endhead
T1 & Demand-modulated model; 00 & GRAY & S13--S18 and older Wu-Xing sources. Coupling priority unresolved; topology and saturation are not new principles.\\
T2 & Arbitrary-$(n,q)$ spectrum; P03--P04 & RED & S07 general circulant theory, S09 ring formulation. Supporting calculation after the Jacobian is derived.\\
T3 & $q=2$ quadratic selection; P05 & ORANGE & S03 Proposition~5 and its two-harmonic dispersion. Need realizable parameter regions and dynamical consequences.\\
T4 & Large-$n$ selected mode; P06 & RED & S03 Proposition~6; fixed-symbol grid maximization is routine. Retain only with useful finite-size/threshold information.\\
T5 & Matched bilinear comparison; P08 & YELLOW & S05 contains MB at five nodes exactly. Potential novelty is the qualified DM--MB comparison, not MB itself.\\
T6 & Explicit Lyapunov reduction; P13 & ORANGE & S02 Appendix~D (including cubic terms), S08 Lemma~1. Known method; model-specific sign information may matter.\\
T7 & Heterogeneous Fourier Green function; P15--P16 & RED & S07 inverse formula and S19 decay theorem. An inverse sum is not a localization theorem.\\
T8 & Robustness bounds; P17--P18 & RED & S10--S11, S20--S21. Generic gap and derivative formulas are known; sharp structured nonlinear bounds remain a separate proposal.\\
\bottomrule
\end{longtable}

\subsection{T1: demand-modulated model --- GRAY}
\begin{description}[style=nextline]
\item[Exact published statement?] Undetermined. No inspected source has
been established as equation-identical to DM; the highest-risk enzymatic
papers lack accessible complete equations. This is a gap, not a novelty finding.
\item[Immediate corollary?] A modeling definition is not a theorem corollary.
It may nevertheless be a reparameterization, restriction or rearrangement of
an existing kinetic family; that comparison is unfinished.
\item[Routine supporting specialization?] The cyclic/five-element topology,
activation/restriction motifs and saturation alone are known ingredients.
\item[Known method, new closed form?] Possibly a new combination of known
kinetic factors; neither provenance nor mathematical value is settled.
\item[Substantive model-specific novelty?] Not established. ``Demand-modulated''
is the frozen descriptive name, not evidence of a new physical law.
\item[Likely referee response?] Ask why these fractions are appropriate and
what consequences they have beyond existing biochemical laws.
\item[Required sharpening.] Complete the S13--S16 equation comparison and
show a specific dynamical contrast under a transparent MB calibration.
Do not rename the term as standard Holling type II or use biomedical
interpretation as mathematical evidence.
\end{description}

\subsection{T2: \texorpdfstring{arbitrary-$(n,q)$}{arbitrary-(n,q)} Fourier spectrum --- RED}
\begin{description}[style=nextline]
\item[Exact published statement?] The general diagonalization theorem is
published (S07). The displayed DM coefficients have not been verified here.
\item[Immediate corollary?] Yes, conditional on establishing that the actual
homogeneous Jacobian is the claimed circulant polynomial in the shift.
\item[Routine supporting specialization?] Yes: spectrum, normality and
the modal Hurwitz test belong in the supporting analysis.
\item[Known method, new closed form?] The substitution may produce a
model-specific expression, not a new spectral method.
\item[Substantive model-specific novelty?] No on the present formulation.
\item[Likely referee response?] Fourier diagonalization of a circulant is
standard; arbitrary offsets do not remove that objection.
\item[Required sharpening.] Use this lemma to support P05/P08/P11. A
genuinely new classification for other $q$ would be separate work and is
not needed in the recommended core.
\end{description}

\subsection{T3: exact quadratic mode selection --- ORANGE}
\begin{description}[style=nextline]
\item[Exact published statement?] Not verified for the DM coefficient
family. S03 already contains a close two-harmonic selection problem.
\item[Immediate corollary?] The basic quadratic maximization is an elementary
consequence of a verified quadratic symbol; the admissible DM parameter
classification and first-instability interpretation are additional tasks.
\item[Routine supporting specialization?] A nearest-cosine rule alone is
routine and should not be presented as a major discovery.
\item[Known method, new closed form?] Yes, potentially. The coefficient
map and its admissible regions may be a new specialization.
\item[Substantive model-specific novelty?] A candidate only if mode regions,
their realizability and a consequential DM--MB contrast are established.
\item[Likely referee response?] Why is optimizing this quadratic materially
different from the Lorenz--96 example?
\item[Required sharpening.] P05 proposal S05 in the decisions file: classify
all admissible maximizing mode pairs, ties and exceptional real modes,
then connect the classification to an actual first-instability path.
Do not infer Hopf from dominance alone.
\end{description}

\subsection{T4: large-network selected-wave-number limit --- RED}
\begin{description}[style=nextline]
\item[Exact published statement?] S03 Proposition~6 establishes an analogous
explicit limit for Lorenz--96; a DM-specific limit has not been certified.
\item[Immediate corollary?] For a fixed continuous symbol with a unique
maximizer, convergence of grid maximizers is generic. Multiple maximizers
and size-dependent parameters require different statements.
\item[Routine supporting specialization?] Yes for the basic ratio limit.
\item[Known method, new closed form?] At most a new coefficient substitution.
\item[Substantive model-specific novelty?] Not the limit alone. Finite-size
instability consequences may be useful as part of the core classification.
\item[Likely referee response?] Dense sampling does not itself justify a
principal asymptotic theorem.
\item[Required sharpening.] P06 proposal S06: quantitative uniform error
for both selected angle and first threshold under explicit nondegeneracy,
with boundaries of applicability identified. Even those estimates need
interpretive value; adding an error order alone does not earn GREEN.
\end{description}

\subsection{T5: exact matched-bilinear comparison --- YELLOW}
\begin{description}[style=nextline]
\item[Exact published statement?] MB at $n=5,q=2$ is exactly contained in
S05. The intended comparison with DM was not found in the inspected corpus;
full priority remains dependent on T1 and the enzymatic sequence.
\item[Immediate corollary?] MB formulas follow familiar cyclic polynomial
analysis. A nonempty region with a specified qualitative contrast is not
supplied merely by the comparator's definition.
\item[Routine supporting specialization?] Separate equilibrium and spectrum
tables would be modest supporting work.
\item[Known method, new closed form?] Yes, potentially a new comparative
specialization. YELLOW is conditional and does not certify originality.
\item[Substantive model-specific novelty?] Possible after a contrast theorem,
not presently demonstrated. Equal parameter count and graph do not ensure
equal equilibrium, response scale or Jacobian signs.
\item[Likely referee response?] Is the contrast robust or an arbitrary
consequence of comparing differently calibrated regimes?
\item[Required sharpening.] P08 proposal S08: freeze the comparison rule
and classify a nonempty open regime with a genuine mode, bifurcation-type,
or criticality distinction; delimit existence of both branches.
\end{description}

\subsection{T6: first Lyapunov coefficient --- ORANGE}
\begin{description}[style=nextline]
\item[Exact published statement?] S02 and S08 publish Fourier reductions;
S02 already includes cubic contributions. No DM coefficient is verified here.
\item[Immediate corollary?] The contraction/selection-rule method is standard
once Taylor tensors, inverses and a simple Hopf setting are justified.
\item[Routine supporting specialization?] A long substituted formula is
usually supporting calculation, even if algebraically laborious.
\item[Known method, new closed form?] This is the most defensible description
of a future DM formula; its closed-form priority is not established.
\item[Substantive model-specific novelty?] A uniform or sharply delimited
criticality classification may be substantive. No sign theorem is supplied.
\item[Likely referee response?] What does the formula decide that the general
Hopf machinery or a single numerical evaluation does not?
\item[Required sharpening.] P13 proposal S13: rigorous sign/nonvanishing
on a meaningful positive parameter region, all selected modes covered by
the claimed range, and explicit exclusions. Do not promise a Bautin point
without its further nondegeneracy conditions.
\end{description}

\subsection{T7: heterogeneous Fourier Green function --- RED}
\begin{description}[style=nextline]
\item[Exact published statement?] The general inverse-circulant identity is
known (S07). A DM forcing numerator may be model-specific.
\item[Immediate corollary?] The linear response follows differentiation
of a nonsingular equilibrium equation and diagonal inversion.
\item[Routine supporting specialization?] Yes; localized forcing does not
automatically imply quantitatively localized response.
\item[Known method, new closed form?] Yes for the coefficient specialization;
that alone has little novelty. S19 supplies a stronger decay baseline.
\item[Substantive model-specific novelty?] Only a sharpened directional,
uniform or critical-scaling result with controlled nonlinear error is a
candidate. It must improve on, not merely restate, generic decay theory.
\item[Likely referee response?] Why is this more than an inverse DFT?
\item[Required sharpening.] P16 proposal S16. Specify graph distance,
parameter family and constants independent of size where claimed.
Distinguish static sensitivity near zero eigenvalues from resonant
harmonic response near a nonzero-frequency Hopf point.
\end{description}

\subsection{T8: robustness bounds and threshold shifts --- RED}
\begin{description}[style=nextline]
\item[Exact published statement?] Generic matrix-gap, resolvent and simple
eigenvalue derivative results are published (S10--S11/S21); the specific
nonlinear parameter certificate is not established here.
\item[Immediate corollary?] The unstructured normal-matrix estimate and a
first-order shift at a simple transverse crossing are routine applications.
\item[Routine supporting specialization?] Yes for a loose sufficient bound
or a left/right eigenvector derivative without further consequences.
\item[Known method, new closed form?] Possibly a useful parameter-level
formula, provided equilibrium displacement and remainder are controlled.
\item[Substantive model-specific novelty?] A sharp certificate for an actual
coefficient uncertainty class is a candidate; no such result is yet verified.
\item[Likely referee response?] Is this just a spectral margin divided by
a rough Lipschitz constant? Is the purported pattern effect nonzero under
the imposed zero-mean constraints?
\item[Required sharpening.] P18 proposal S18: explicit structured parameter
radius with positivity/local existence and equilibrium motion included;
quantify conservatism or sharp scaling. Compare any disorder-stabilization
or second-order mode-mixing claim against S24. Do not identify a general
full-block stability radius with a repeated/diagonal coefficient radius.
\end{description}
